\section{BFS Growth as MCMC Warm-Start}
\label{sec:warmstart}

\subsection{Motivation}

MCMC redistricting chains---\recom{}, Forest ReCom, and their variants---
require an initial feasible plan from which the chain begins.
The quality of this warm-start affects two aspects of chain behaviour:
\begin{enumerate}
\item \textbf{Burn-in length.} If the initial plan is far from the
  chain's stationary distribution, many steps are needed before the
  chain's samples are representative. This means results from early
  steps must be discarded as burn-in.
\item \textbf{Autocorrelation.} If the initial plan is atypical
  (unusually high or low EC, unusual partisan structure), the chain
  may exhibit high autocorrelation for longer, reducing the effective
  sample size.
\end{enumerate}

The default warm-start in the B-series platform is a single \metis{} call.
\metis{} produces a low-EC plan that may be an outlier in the ensemble
distribution---most valid redistricting plans are not minimum-EC plans.
\bfsg{} produces a higher-EC, more geographically typical plan that may
lie closer to the centre of the ensemble distribution.

\subsection{Theoretical Argument}

Let $\pi$ denote the stationary distribution of the \recom{} chain over
valid redistricting plans of a given state.
The typical plan under $\pi$ is characterised by the ensemble's modal
EC and PP values, which are substantially higher than the METIS minimum
and somewhat below the CVD or BFS Growth values.

If the initial plan $x_0$ has EC close to the ensemble mode, the chain
needs fewer steps to mix because each ReCom merge-split proposal is more
likely to produce a plan similar to $x_0$---and similarly valued by $\pi$.
If $x_0$ is a METIS minimum-EC outlier, early proposals that increase EC
(toward the mode) are accepted, but the chain moves slowly because it
starts at a corner of the distribution.

\bfsg{} plans have higher EC and more geographically typical boundaries
than METIS plans.
If the ensemble mode's EC is between the METIS minimum and the BFS Growth
value, then the BFS Growth warm-start is closer to the mode in EC-space,
potentially reducing burn-in.

\subsection{Expected Empirical Results}

We characterise the expected experimental findings for future validation:

\paragraph{NC ($k=14$).}
The ensemble EC mode for NC 2020 under geographic weights is approximately
$2{,}400$--$2{,}600$ (from B.7 500-seed sweep data~\citep{b7paper}).
\metis{} single-seed EC is approximately $2{,}891$ (Table~\ref{tab:main-comparison}),
which is \emph{above} the mode---suggesting that a single METIS call may
not be the lowest-EC plan.
\bfsg{} EC is $3{,}847$, above the METIS baseline and substantially above
the ensemble mode.

In this case, neither warm-start is at the mode.
The METIS warm-start is closer to the mode in EC-space.
Expected finding: \metis{} warm-start has shorter burn-in on NC.

\paragraph{WI ($k=8$).}
WI's smaller graph ($1{,}406$ tracts) means the ensemble distribution is
tighter.
Both METIS and BFS Growth warm-starts are expected to burn in quickly.
Expected finding: comparable burn-in lengths for both warm-starts.

\paragraph{TX ($k=38$).}
TX's large graph and diverse geography create a wide ensemble distribution.
The expected EC mode is in a range where both METIS and BFS Growth plans
are outliers.
Expected finding: both warm-starts require substantial burn-in; BFS Growth
may provide faster autocorrelation decay due to more geographically
regular boundaries that are less disruptive to the chain's spanning-tree
proposals.

\subsection{Autocorrelation Decay}

Let $f: \mathcal{P} \to \mathbb{R}$ be a plan statistic (e.g., EC, mean PP,
partisan seat count).
The autocorrelation at lag $\ell$ is:
\[
  \rho_f(\ell) = \frac{\mathrm{Cov}(f(x_t),\; f(x_{t+\ell}))}
                      {\mathrm{Var}(f(x_t))},
\]
where $(x_t)$ is the Markov chain.
The \emph{integrated autocorrelation time} $\tau_f = \sum_{\ell=1}^\infty \rho_f(\ell)$
quantifies how many steps are needed for effectively independent samples.

A smaller $\tau_f$ from a BFS Growth warm-start relative to a METIS
warm-start would indicate faster mixing when starting from the BFS plan.
The key hypothesis is:

\begin{quote}
\textit{For states where the BFS Growth plan's EC is closer to the
ensemble mode than the METIS plan's EC, the ReCom chain warm-started
from BFS Growth has smaller $\tau_f$ for EC-correlated statistics.}
\end{quote}

This hypothesis is testable via paired chain runs on NC, WI, and TX
with 10{,}000 steps and 50 burn-in discarded.
The result is left as future empirical work; the B.23 theoretical
framework and expected directions are stated here.

\subsection{Practical Recommendation}

Given the current analysis, the practical warm-start recommendation is:

\begin{itemize}
\item \textbf{Default: \metis{} warm-start.}
  \metis{} is better understood, and its EC properties are well-documented.
  Use \metis{} as the default warm-start until BFS Growth warm-starts
  are empirically validated.

\item \textbf{Alternative: \bfsg{} warm-start for geographic analysis.}
  When the research question concerns geographic compactness or the
  spatial structure of valid plans, a BFS Growth warm-start provides a
  more geographically typical initial plan.
  The BFS Growth plan's higher EC means the chain starts with more
  boundary-swap flexibility, potentially exploring a wider range of
  plans early.

\item \textbf{Multi-start: both.}
  For rigorous ensemble analysis, run two parallel chains from both
  warm-starts and compare their convergence diagnostics.
  Agreement between the two chains' stationary-phase distributions
  confirms mixing; disagreement signals that longer chains are needed.
\end{itemize}

\subsection{Implementation}

To use \bfsg{} as a warm-start for a \recom{} ensemble:

\begin{verbatim}
# Step 1: Generate BFS Growth initial plan
redist state --state NC --year 2020 \
  --structure bfs-growth --search single

# Step 2: Use the output plan as ReCom warm-start
redist ensemble --state NC --year 2020 \
  --initial-plan runs/nc_bfs/2020/nc_plan.json \
  --chain recom --steps 10000
\end{verbatim}

The \texttt{--initial-plan} flag accepts any previously generated plan
in the platform's JSON format, including BFS Growth plans.
The audit chain (Section~3 of the spec) records the \texttt{base\_seed}
and \texttt{bfs\_seed} in \texttt{runs/{label}/{year}/index.json},
enabling full reproducibility of the warm-start.
